\chapter{Veri Setleri ve Değerlendirme}
\label{ch:veri-setleri}

\section{Neden Bu Bölüm Var}
\label{sec:neden-veri-seti}

Raporun buraya kadarki doğruluk sayılarının tamamı bir \emph{parite}
ölçümüdür: fp16 motorlarının çıktısı, aynı modelin fp32 çıktısına ne
kadar yakın kaldı (\cref{sec:olcum-yontemi}). Bu, motorun doğru
derlendiğini kanıtlar; modelin \emph{ne kadar iyi takip ettiğini}
kanıtlamaz, çünkü karşılaştırılan referans modelin kendi çıktısıdır.
Dışarıdan bir doğruluk sayısı için iki şey gerekir: bir veri seti ve o
veri setinde kare başına bir \textbf{yer doğrusu} (ground truth, GT).

Projede kullanılan Anti-UAV410~\cite{huang2023antiuav410} bu boşluğu
yalnızca yarısına kadar doldurur: etiketleri \emph{kutu}dur. Kutu, IoU
ve başarı eğrisi verir (\texttt{src/accuracy.py}), ama EdgeTAM'ın
ürettiği şey bir \emph{maske}dir ve maskenin kalitesi kutuyla
ölçülemez. Maske kalitesini ölçen standart metrik J\&F'tir
(\cref{sec:jf-metrigi}) ve J\&F, kare başına \textbf{yoğun (piksel
düzeyinde) maske GT'si} ister.

Bu bölüm, aday veri setlerini tam olarak bu ölçütle eler:

\begin{enumerate}
  \item \textbf{Gerçekten indirilebilir mi?} Yayınlanmış bir bağlantı
        değil, çalışan bir dağıtım (figshare / Hugging Face / Drive /
        kurum arşivi) ve açık bir lisans.
  \item \textbf{Yoğun maske GT'si var mı?} Kutu, GPS ya da sınıf
        etiketi J\&F üretmez.
  \item \textbf{Havadan mı?} Tabloda \hava\ ile işaretli.
  \item \textbf{Video mu?} EdgeTAM'ın ölçtüğü şey hafıza mekanizmasıdır;
        tek karelik bir set onu ne eğitir ne ölçer.
\end{enumerate}

\section{Aday Veri Setleri}
\label{sec:aday-veri-setleri}

\begin{table}[H]
  \centering
  \caption{İncelenen veri setleri. \hava\ = havadan (UAV) çekim.
    ``Yoğun maske'' = piksel düzeyinde semantik/örnek maskesi;
    parantez içindeki aralık, videolarda hangi karelerin etiketlendiğini
    gösterir. HF = Hugging Face. $\dagger$ = kendi çekimi değil, başka
    veri setlerinden türetilmiş (\cref{sec:ortak-kaynaklar}).
    Satırlar öncelik sırasına göre gruplanmıştır.}
  \label{tab:veri-setleri}
  \footnotesize
  \setlength{\tabcolsep}{3pt}
  \renewcommand{\arraystretch}{1.18}
  \begin{tabular}{@{}
      >{\raggedright\arraybackslash}p{21mm}
      >{\raggedright\arraybackslash}p{15mm}
      >{\raggedright\arraybackslash}p{14mm}
      >{\raggedright\arraybackslash}p{18mm}
      >{\raggedright\arraybackslash}p{10mm}
      >{\raggedright\arraybackslash}p{24mm}
      >{\raggedright\arraybackslash}p{25mm}
      >{\raggedright\arraybackslash}p{14mm}@{}}
    \toprule
    \textbf{Veri seti} & \textbf{Yıl · Yayın} & \textbf{Modalite} &
    \textbf{Çözünürlük} & \textbf{Tip} & \textbf{Yer doğrusu (GT)} &
    \textbf{Hacim} & \textbf{Erişim} \\
    \midrule
    \multicolumn{8}{@{}l}{\textit{\color{grimetin} Mevcut referans --- bu projede zaten kullanılıyor}}\\
    Anti-UAV410\hava & 2023 · TPAMI & TIR & 640×512 & Video &
      Kutu + görünürlük bayrağı & 410 video, 438\,K kutu & Drive \\
    \midrule
    \multicolumn{8}{@{}l}{\textit{\color{yesilvurgu} Öncelikli --- yoğun maske GT'si var, indirmesi doğrulandı}}\\
    Kust4K\hava & 2025 · Sci.\ Data & RGB+TIR (hizalı) & 640×512 &
      Statik çift & Yoğun maske, 8 sınıf & 4\,024 hizalı çift & figshare \\
    MVUAV\hava & 2024 · NeurIPS & RGB+TIR & 1920×1080 & Video &
      Yoğun maske, 36 sınıf (25 karede bir) &
      413 video / 53\,828 kare çifti; 2\,183 etiketli & proje sayfası \\
    Caltech Aerial RGB-T\hava & 2024 · ECCV & RGB+TIR & 960×600 /
      640×512 & Video & Yoğun maske, 10 sınıf &
      37 çekim (18'i havadan), 4\,195 maske & Caltech arşivi \\
    AeroVIS\hava\turetilmis & 2026 · arXiv & RGB & kaynak setlere göre değişken &
      Video & \textbf{Örnek maskesi + kimlik izi}, 9 sınıf &
      117 video / 49\,204 kare, 8\,279 iz & Drive (MIT) \\
    SegFly\hava & 2026 · ECCV & RGB+TIR & 5472×3648 / 640×512 & Statik &
      Yoğun maske, 15 sınıf & 20\,606 RGB + 15\,007 RGB-T çifti & HF \\
    MVSeg\turetilmis & 2023 · CVPR & RGB+TIR & karışık, $\leq$1280×720 & Video &
      Yoğun maske, 26 sınıf (15 karede bir) &
      738 video, 3\,545 maske & proje sayfası \\
    \midrule
    \multicolumn{8}{@{}l}{\textit{\color{grimetin} İkincil --- GT ya kutu, ya sentetik, ya da hacim/çözünürlük yetersiz}}\\
    VTUAV\hava & 2022 · CVPR & RGB+TIR & 1920×1080 & Video &
      Kutu (500 dizi) + maske (100 video) &
      500 dizi, $\sim$1,7\,M kare çifti & Drive / Baidu \\
    VDD\hava\turetilmis & 2023 · arXiv & RGB & yüksek çöz. & Statik &
      Yoğun maske, 7 sınıf & 400 görüntü & HF \\
    FlyAwareV2\hava\turetilmis & 2026 · SPIC & RGB + derinlik & --- & Statik &
      Yoğun maske + derinlik (sentetik + gerçek) &
      100\,K+ örnek, $\sim$296\,GB & resmî site, \mbox{GPL-3.0} \\
    Multi-View UAV\hava & 2025 · arXiv & RGB + derinlik & 400×300 &
      Statik (uçuş sıralı) & Yoğun maske + derinlik (CARLA) &
      357\,690 kare (5 kamera), 140\,GB & HF (MIT) \\
    \bottomrule
  \end{tabular}
\end{table}

\section{Veri Setleri Hakkında Notlar}
\label{sec:veri-seti-notlari}

\subsection*{Öncelikli}

\paragraph{Kust4K~\cite{kust4k2025}.}
Bu listedeki tek veri seti, projenin girdi formatıyla \emph{birebir}
aynı: 640×512 termal, RGB'siyle piksel düzeyinde hizalanmış. Aynı
çözünürlük, \texttt{crop512} modunun sıfır yeniden boyutlandırmayla
çalıştığı çözünürlüktür (\cref{sec:onisleme-analizi}) --- yani ön-işleme
maliyeti ölçümü bozmaz. 4\,024 çift, 8 sınıf (yol, bina, motosiklet,
araba, kamyon, ağaç, insan, trafik ekipmanı); 2\,514'ü gündüz, 857'si
gece, 653'ü düşük ışıklı gece.
\textbf{Zayıf yanı:} bağımsız kareler, video değil; hafıza mekanizmasını
ölçemez, ama tek kare maske kalitesi ve termal alan uyumu için en temiz
kaynak.

\paragraph{MVUAV~\cite{ji2024mvuav}.}
Havadan, eğik kuşbakışı RGB-T \emph{video}: 413 video, 53\,828 kare
çifti, bunların 2\,183'ü 36 sınıfta piksel düzeyinde etiketli (25
karede bir anahtar kare). Gündüz/gece, alacakaranlık ve zifiri karanlık
sahneler içerir --- yani \cref{sec:hafiza-toparlanma}'deki hafıza
kirlenmesi senaryosunun (ani kontrast değişimi) doğal olarak bulunduğu
bir set. \textbf{Zayıf yanı:} etiketler zamanda seyrek; J\&F yalnızca
anahtar karelerde hesaplanabilir, ki bu zaten MVSS literatürünün
standardıdır.

\paragraph{Caltech Aerial RGB-T~\cite{lee2024caltechrgbt}.}
Termal kamerası \textbf{FLIR ADK, 640×512, 60\,Hz} --- yine projenin
native çözünürlüğü. 37 çekimin 18'i havadan (çoğunlukla 40\,m irtifa),
toplam 4\,195 yoğun maske, 10 sınıf, doğal ortamlar (göl, nehir, kıyı,
çöl, orman). GPS ve IMU da senkron kayıtlı. \textbf{Güçlü yanı:} donanım
senkronize RGB-T çiftleri + gerçek uçuş hareketi; CC BY-SA lisansıyla
doğrudan indirilebiliyor. \textbf{Zayıf yanı:} sınıflar doğa odaklı,
drone/araç gibi küçük hedef yok.

\paragraph{AeroVIS~\cite{aerovis2026}.}
\textbf{Bu listede EdgeTAM'ın görevine en yakın veri seti.} Diğerleri
semantik segmentasyon (piksel $\rightarrow$ sınıf); AeroVIS ise
\emph{örnek} segmentasyonu \emph{artı kimlik takibi} (piksel
$\rightarrow$ nesne, nesne video boyunca aynı kimlikle) --- yani tam
olarak EdgeTAM'ın ürettiği çıktı tipi. 117 video, 49\,204 kare, 9 UAV
nesne kategorisi, 8\,279 geçerli iz, YTVIS biçiminde. VisDrone, UAVDT ve
SeaDronesSee'nin elle çizilmiş kutuları SAM3 ile maskeye çevrilip elden
geçirilerek üretilmiş. \textbf{Zayıf yanı:} maskeler otomatik üretim
kaynaklı (elle denetlenmiş ama elle çizilmemiş) ve RGB; termal yok.
Kendi çekimi olmadığı için kaynak setleriyle çakışmaya dikkat
(\cref{sec:ortak-kaynaklar}).

\paragraph{SegFly~\cite{segfly2026}.}
Ölçek olarak listenin en büyüğü: 20\,606 RGB görüntü ve 15\,007 hizalı
RGB-T çifti, 15 sınıf, 30/40/50\,m irtifa. Termal yine 640×512.
Etiketlerin yalnızca \%2,84'ü elle çizilmiş; kalanı 2D--3D--2D geometrik
yayılımla üretilmiş --- bu yüzden hacim büyük ama etiket gürültüsü
sıfır değil. \textbf{Zayıf yanı:} bağımsız hava fotoğrafları (video
değil) ve CC BY-NC-SA (ticari kullanım yok).

\paragraph{MVSeg~\cite{ji2023mvseg}.}
MVUAV'ın yer seviyesi kardeşi ve öncülü: 738 RGB-T video, 3\,545 maske,
26 sınıf, 15 karede bir etiket. Havadan \emph{değil} --- tabloda yıldızsız
olmasının nedeni bu. Buna rağmen listede kalmasının nedeni pratik: MVUAV
ile aynı ekipten, aynı biçimde ve aynı değerlendirme koduyla geliyor, yani
ikisi tek bir işle birlikte kullanılabilir. OSU, INO, KAIST ve RGBT234'ten
derlendiği için çözünürlük karışıktır ve setin tamamı
türetilmedir (\cref{sec:ortak-kaynaklar}).

\subsection*{İkincil}

\paragraph{VTUAV~\cite{zhang2022vtuav}.}
Hacim olarak listenin devi: 500 dizi, $\sim$1,7 milyon hizalı 1920×1080
RGB-T kare çifti. Ama ana etiketi \textbf{kutu}dur; maske yalnızca 100
videoluk bir alt küme için var. Yani doğruluk ölçümü için değil,
\emph{hız ve dayanıklılık} ölçümü için değerli: uzun diziler,
\cref{sec:hafiza-toparlanma}'deki hafıza kirlenmesinin uzun videoda
birikip birikmediğini görmek için doğal bir test yatağı.

\paragraph{VDD~\cite{cai2023vdd}.}
Yoğun etiketli, temiz, açık --- ama 400 görüntü ve yalnızca RGB. Bu
hacim bir fine-tune için az, bir benchmark için sınırda. UDD ve UAVid'in
yeniden etiketlenmiş halleriyle (IDD paketi) birlikte geldiği için asıl
değeri tek başına değil, o üçlüyü tek etiket standardında toplaması ---
ki bu aynı zamanda çakışma kaynağıdır (\cref{sec:ortak-kaynaklar}).

\paragraph{FlyAwareV2~\cite{flyawarev2}.}
100\,K'yı aşan örnek, dört hava koşulu (güneşli, yağmurlu, sisli, gece),
üç irtifa (20/50/120\,m), RGB + derinlik + semantik. \textbf{Sorun:}
büyük kısmı CARLA ile üretilmiş sentetik veri, gerçek kısmı ise yalnızca
$\sim$2\,K kare: VisDrone (eğitim) ve UAVid (test), hava koşulu
artırımıyla çoğaltılmış. Tam sürüm $\sim$296\,GB, termal yok. Bu iki
kaynak da başka setlerde geçiyor (\cref{sec:ortak-kaynaklar}).
Alan uyumu (domain gap) açısından, projenin gerçek termal kayıtlarına
sentetik RGB'den daha uzak bir veri kaynağı.

\paragraph{Multi-View UAV~\cite{multiviewuav2025}.}
357\,690 kare kulağa çok geliyor ama kareler 400×300 ve tamamı CARLA
simülasyonundan; her ``kare'' beş kameralı bir kurulumun tek görüntüsü.
Semantik maske ve derinlik hazır, lisans MIT, ama 400×300 EdgeTAM'ın
512 girdi çözünürlüğünün altında --- yani her karede yukarı ölçeklemek
gerekir, bu da ölçülen maske kalitesini veri setinin kendi
çözünürlüğüyle sınırlar.

\section{Ortak Kaynaklar ve Çakışma Riski}
\label{sec:ortak-kaynaklar}

\Cref{tab:veri-setleri}'ndeki on setin dördü kendi çekimi değildir:
daha eski, halka açık veri setlerinin yeniden etiketlenmiş ya da
yeniden paketlenmiş halleridir. Bu, tek başına bir kusur değil ---
kutuyu maskeye çevirmek gerçek bir katkıdır --- ama iki set aynı ham
kareleri paylaşıyorsa, birinde eğitip diğerinde J\&F raporlamak
\textbf{kirlenmiş bir ölçüm} verir. Kaynaklar:

\begin{table}[H]
  \centering
  \caption{Türetilmiş veri setleri ve kaynakları. Tabloda $\dagger$ ile
    işaretli satırlar bunlardır; diğer altı set kendi çekimidir
    (Multi-View UAV'ın ``çekimi'' CARLA simülasyonudur).}
  \label{tab:veri-seti-kaynaklari}
  \footnotesize
  \setlength{\tabcolsep}{4pt}
  \renewcommand{\arraystretch}{1.18}
  \begin{tabular}{@{}
      >{\raggedright\arraybackslash}p{24mm}
      >{\raggedright\arraybackslash}p{40mm}
      >{\raggedright\arraybackslash}p{68mm}@{}}
    \toprule
    \textbf{Set} & \textbf{Kaynak} & \textbf{Ne yapılmış} \\
    \midrule
    AeroVIS & VisDrone, UAVDT, SeaDronesSee &
      Tamamı türetilmiş. Kaynakların elle çizilmiş kutuları konum ve
      kimlik önceli olarak kullanılıp SAM3 ile maskeye çevrilmiş,
      sonra elden geçirilmiş. \\
    FlyAwareV2 & gerçek: VisDrone (eğitim) + UAVid (test);
      sentetik: CARLA / SynDrone &
      $\sim$2\,K gerçek kare, hava koşulu artırımıyla çoğaltılmış;
      $\sim$288\,K kare simülasyondan. \\
    VDD & (kendi 400 görüntüsü) + UDD, UAVid &
      VDD'nin 400 görüntüsü özgün; birlikte dağıtılan IDD paketi UDD ve
      UAVid'i VDD etiket standardında yeniden etiketliyor. \\
    MVSeg & OSU, INO, KAIST, RGBT234 &
      Mevcut RGB-T video kaynakları MVSS görevine uyarlanıp piksel
      düzeyinde yeniden etiketlenmiş. \\
    \bottomrule
  \end{tabular}
\end{table}

\subsection*{Hangi İkililer Çakışıyor}

\begin{itemize}
  \item \textbf{VisDrone $\rightarrow$ AeroVIS + FlyAwareV2.} Listedeki
        tek gerçek kare düzeyinde çakışma riski budur. FlyAwareV2'nin
        gerçek \emph{eğitim} bölümü VisDrone'dur ve AeroVIS'in videoları
        da kısmen VisDrone'dan gelir. Bu ikisini bir eğitim/değerlendirme
        çifti olarak kullanmak güvenli değildir.
  \item \textbf{UAVid $\rightarrow$ FlyAwareV2 + VDD/IDD.} FlyAwareV2'nin
        \emph{test} bölümü UAVid'dir; VDD ile gelen IDD paketi de UAVid'i
        yeniden etiketler. IDD'de eğitip FlyAwareV2 testinde ölçmek aynı
        kareleri iki kez görmek olur.
  \item \textbf{CARLA $\rightarrow$ FlyAwareV2 + Multi-View UAV.} Burada
        paylaşılan şey kareler değil, \emph{simülatör}. Sahneler ve
        kameralar farklı, yani veri sızıntısı yok; ama ikisi de aynı
        render motorunun görsel imzasını taşır, dolayısıyla ikisinde de
        iyi çalışan bir model gerçek termal kayıtta aynı sonucu vermez.
  \item \textbf{RGBT234 $\rightarrow$ MVSeg.} RGBT234 aslında bir RGB-T
        \emph{takip} benchmark'ıdır. Bugün için zararsız, ama ileride
        takip doğruluğu RGBT234 üzerinde ölçülürse MVSeg ile örtüşür.
\end{itemize}

\subsection*{Sonuç: Öncelikli Blok Kendi İçinde Temiz}

Kritik nokta şudur: öncelikli bloktaki tek türetilmiş set
\textbf{AeroVIS}'tir ve onun kaynakları (VisDrone, UAVDT, SeaDronesSee)
bloktaki diğer beş setin --- Kust4K, MVUAV, Caltech Aerial RGB-T,
SegFly, MVSeg --- hiçbirine değmez. Yani
\cref{tab:veri-setleri}'nin öncelikli altısı \textbf{birlikte
kullanılabilir}; aralarında paylaşılan tek bir kare yoktur.

Çakışma, ikincil bloğa geçildiğinde başlar. Pratik kural: FlyAwareV2 ve
VDD/IDD ile eğitim yapılacaksa, AeroVIS artık tarafsız bir
değerlendirme kümesi değildir. En temiz kurulum, hem eğitimi hem
değerlendirmeyi \textbf{öncelikli blokla sınırlamaktır}: örneğin
Kust4K ve MVUAV'da eğitip J\&F'i AeroVIS'te ölçmek, aralarında tek bir
paylaşılan kare olmadığı için güvenlidir.

\section{Elenen Bağlantılar}
\label{sec:elenen-veri-setleri}

Üç aday, yukarıdaki ölçütlerin ilk ikisinden geçemedi:

\begin{itemize}
  \item \textbf{UAV-VisLoc}~\cite{uavvisloc2024} --- 6\,742 havadan
        görüntü, 11 uydu haritası, 16,4\,GB. Açık ve indirilebilir, ama
        yer doğrusu \emph{GPS koordinatı}; maske ya da kutu yok. Görev
        görsel konumlandırma, segmentasyon değil.
  \item \textbf{CPD-UAV}~\cite{cpduav2026} --- 1\,061 yüksek çözünürlüklü
        görüntü, ortamla görsel olarak kaynaşmış insanların tespiti.
        Etiket yalnızca kutu, hacim küçük, video yok.
  \item \textbf{Scilit kaydı} (\texttt{d95b5f3c\ldots}) --- sayfa
        erişime kapalı (HTTP 403) ve kayıt başka bir yerde
        eşleştirilemedi; hangi yayına işaret ettiği doğrulanamadığı için
        listeden çıkarıldı.
\end{itemize}

Ayrıca \textbf{Kust4K'nın figshare bağlantısı ile Nature Scientific
Data makalesi aynı veri setidir}; ham listede iki ayrı satır olarak
görünüyorlardı, tabloda tek satıra indirildi.

\section{J\&F: Video Nesne Segmentasyonunun Standart Metriği}
\label{sec:jf-metrigi}

Yukarıdaki elemenin tamamı tek bir metriğin gereksinimlerinden çıkıyor,
o yüzden metriğin kendisi burada tanımlanıyor. J\&F, DAVIS
benchmark'ıyla~\cite{perazzi2016davis,ponttuset2017davis} birlikte
gelmiş ve video nesne segmentasyonunun fiilî standardı olmuştur.
Bu ailenin tamamı --- SAM~2, EdgeTAM, SAMURAI, SAM2Long ---
doğruluğunu bu sayıyla raporlar; yani bu projenin sonuçlarını
literatürle karşılaştırılabilir kılacak sayı da budur.

İki yarımın ortalamasıdır. $M$ tahmin edilen maske, $G$ yer doğrusu
maskesi olmak üzere:

\paragraph{$\mathcal{J}$ --- bölge benzerliği (region similarity).}
Maskelerin Jaccard indeksi, yani maske IoU'su:
\begin{equation}
  \mathcal{J} = \frac{|M \cap G|}{|M \cup G|}
  \label{eq:jaccard}
\end{equation}
``Doğru pikselleri buldun mu?'' sorusunu ölçer.

\paragraph{$\mathcal{F}$ --- kontur doğruluğu (contour accuracy).}
Yalnızca maskelerin \emph{sınır} piksellerine bakar. Tahminin sınır
noktaları, yer doğrusunun sınır noktalarıyla bir $\tau$ toleransı
içinde eşleştirilir; eşleşenlerin oranı kesinliği $P_c$ ve
duyarlılığı $R_c$ verir:
\begin{equation}
  \mathcal{F} = \frac{2\,P_c\,R_c}{P_c + R_c}
  \label{eq:fmeasure}
\end{equation}
``Nesnenin kenarını doğru yerde çizdin mi?'' sorusunu ölçer. DAVIS'te
tolerans görüntü köşegeninin \%0,8'idir (1920×1080 için
$\approx$18~piksel).

\paragraph{Birleşim.}
\begin{equation}
  \mathcal{J\&F} = \frac{\mathcal{J} + \mathcal{F}}{2}
  \label{eq:jf}
\end{equation}
Kare başına hesaplanır ve klip boyunca ortalanır; komutun verildiği ilk
kare hesaba katılmaz, çünkü orada maske zaten kullanıcıdan gelmiştir.

\subsection*{Neden İki Yarım? Kısa Bir Örnek}

$\mathcal{J}$ bir \emph{alan} ölçüsüdür ve nesne büyüklüğüne karşı
tarafsız değildir. Aynı mutlak hata, nesne küçüldükçe çok daha ağır
cezalandırılır. Tahminin yer doğrusuyla aynı şekilde ama 5~piksel kaymış
olduğu iki durum:

\begin{table}[H]
  \centering
  \caption{Aynı 5 piksellik kayma, iki farklı nesne boyutunda.}
  \label{tab:jf-ornek}
  \small
  \begin{tabular}{lrrrr}
    \toprule
    Nesne & Alan (px) & Kesişim & Birleşim & $\mathcal{J}$ \\
    \midrule
    Büyük kare, 100×100 & 10\,000 & 9\,500 & 10\,500 & \textbf{0,905} \\
    Küçük kare, 20×20   &    400  &   300  &    500  & \textbf{0,600} \\
    \bottomrule
  \end{tabular}
\end{table}

Hata birebir aynıdır; $\mathcal{J}$ 0,905'ten 0,600'e düşer. Bu,
proje için teorik bir ayrıntı değil: uzaktaki bir drone karede birkaç
yüz piksellik bir hedeftir, yani $\mathcal{J}$'nin acımasız olduğu
rejimin tam içindedir.

$\mathcal{F}$ ise ters yönde bir körlüğü kapatır. Tahminin, yer
doğrusunun 5~piksel \emph{şişmiş} hali olduğunu düşünelim: nesne doğru
yerde, doğru şekilde, sadece kenarları taşmış. $\mathcal{J}$ hâlâ
yüksek çıkar (50~piksel yarıçaplı bir daire için $\approx$0,83), ama her
sınır noktası olması gerektiği yerden tam 5~piksel uzaktadır ---
tolerans 5'in altındaysa $\mathcal{F}=0$, üstündeyse $\mathcal{F}=1$.
Yani $\mathcal{F}$, ``maskeyi kabaca doğru yere koydum ama kenarı
dağınık'' hatasını yakalar; $\mathcal{J}$ bunu neredeyse hiç görmez.
Bu hata, düşük çözünürlüklü maske logitlerinin kaynak çözünürlüğe geri
ölçeklenmesinden (\texttt{post} aşaması,
\cref{tab:kare-bilesenleri}) doğan tipik hatadır.

\subsection*{Bu Projede Ne Anlama Geliyor}

\begin{itemize}
  \item Raporda geçen \textbf{0,9993 ortalama maske IoU} tam olarak
        \cref{eq:jaccard}'teki $\mathcal{J}$'dir --- ama iki
        \emph{kendi} çıktımız (fp16 ve fp32) arasında. Bir parite
        ölçümüdür; J\&F değildir ve bir doğruluk sayısı olarak
        okunmamalıdır.
  \item Gerçek bir J\&F için dış bir yoğun maske GT'si şarttır.
        Anti-UAV410'un kutu etiketleri bunu veremez;
        \cref{tab:veri-setleri}'nin ``öncelikli'' bloğu tam da bu
        eksiği kapatmak için seçilmiştir.
  \item Küçük hedeflerde $\mathcal{J}$ ile $\mathcal{F}$ ters yönlere
        gider: nesne küçüldükçe $\mathcal{J}$ sertleşir, sabit
        toleranslı $\mathcal{F}$ ise görece cömertleşir. Bu yüzden
        drone takibinde ortalamayı tek başına raporlamak yanıltıcıdır;
        \textbf{iki yarımın ayrı ayrı verilmesi} daha bilgilendiricidir.
\end{itemize}
